\begin{tcolorbox}[colback=red!3!white,colframe=red!50!black,boxrule=0.8pt,title={Korrekturhinweise},fonttitle=\bfseries]
\textbf{Korrekturen aus der Rechenprüfung (30.~Sept.~2026).} Die folgenden Stellen sind im Text korrigiert bzw. vermerkt; jede trägt dort einen datierten Hinweis mit dem vorherigen Wert:
\begin{itemize}
\item Äquivalenztabelle Tau: $r_\tau = 25/9$, Genauigkeit $99.99 \to 99.63\,\%$ (die $99.99\,\%$ gehörten zum gefitteten $2.7675$).
\end{itemize}
\end{tcolorbox}

% Chapter file: 046_Teilchenmassen_De_ch.tex
% Source: 046_Teilchenmassen_De.tex
% German version for T0 Theory Book
% FULLY CORRECTED for KOMA-Script and unicode-math compatibility
% ALL \rm commands replaced with proper LaTeX2e/KOMA-Script syntax

\begin{tcolorbox}[colback=red!3!white,colframe=red!50!black,boxrule=0.8pt,title={Korrekturhinweise},fonttitle=\bfseries]
\textbf{Korrekturen aus der Rechenprüfung (30.~Sept.~2026).} Die folgenden Stellen sind im Text korrigiert bzw. vermerkt; jede trägt dort einen datierten Hinweis mit dem vorherigen Wert:
\begin{itemize}
\item Neutrinomassen $9.1/1.9/18.0$~meV: Vermerk -- dieselben Werte wie im Neutrinoteil von Dok.~006 (R109, nicht tragfähig), ohne gemeinsamen Anker, mit den Oszillationsdaten unverträglich; maßgeblich ist Dok.~340.
\item Summe der Neutrinomassen (Tabelle): $29.8 \to 29.0$~meV.
\item Geladenes Lepton der 4.~Generation bei $5.7$~GeV: Vermerk -- Rechnung richtig, experimentell ausgeschlossen (LEP, etwa $100$~GeV).
\item Direkte Methode $E = 1/\xi_e$, $1/\xi_\mu$: Vermerk -- reine Zahlen $7500$ bzw. $2343.75$, kein Schritt zu MeV (vgl. R109). Tau: $1777.1 \to 1783.4$~MeV ($99.99 \to 99.63\,\%$); Elektron mit $r_e = 4/3$: $0.511 \to 0.505$~MeV ($99.98 \to 98.82\,\%$); Strange: $95.0 \to 94.76$~MeV (nicht exakt); Charm: $1.279 \to 1.284$~GeV; Quark-Referenzmassen auf PDG~2024; Durchschnitt $99.6 \to 98.7\,\%$ (auch Abstract und Status). Neutrino-Exponenten $3, 25/9 \to 2, 5/3$; Äquivalenztabelle: $m_\nu$ in GeV $10^{-6} \to 10^{-12}$, $f_i$ neu berechnet.
\end{itemize}
\end{tcolorbox}

	\section*{Abstract}
	Das T0-Modell bietet zwei mathematisch äquivalente, aber konzeptionell verschiedene Berechnungsmethoden für Teilchenmassen: Die direkte geometrische Methode und die erweiterte Yukawa-Methode. Beide Ansätze verwenden als geometrische Konstante nur $\xipar = \frac{4}{3} \times 10^{-4}$; die rationalen Koeffizienten $r_i$ werden in der Etablierungsphase aus gemessenen Massen festgelegt. Die Dokumentation enthält zudem einen Ansatz für die Neutrino-Quantenzahlen sowie eine quantenfeldtheoretische Herleitung der $\xi$-Konstante durch EFT-Matching und 1-Loop-Rechnungen. Die systematische Behandlung aller Teilchen, einschließlich der Neutrinos mit der angenommenen doppelten $\xi$-Unterdrückung, zeigt die universelle Struktur des Ansatzes. Die durchschnittliche Abweichung von etwa 1.3\% über die geladenen Fermionen \textit{(Korrigiert am 30.~Sept.~2026: vorher \enquote{weniger als 1\%} -- Mittelwert der Genauigkeiten in Tabelle~\ref{tab:complete_validation} mit den tabellierten $r_i$ und PDG-2024-Referenzmassen: $98.7\,\%$.)} bedeutet, verglichen mit über zwanzig freien Standardmodell-Parametern, eine deutliche Reduktion der Eingangsgrößen.

	\medskip\noindent\textbf{Hinweis zur Fassung.} Die deutsche Fassung dieses Dokuments ist umfangreicher als die englische. Bei früheren Überarbeitungen wurden in der englischen Fassung Teile gekürzt oder anders gefasst. Bei Abweichungen gilt die umfangreichere deutsche Fassung.


	\medskip\noindent\textbf{Statusmarker.} Aussagen mit \textbf{[S]} sind \emph{spekulativ}: ein Hinweis oder eine Vermutung, die im Korpus noch nicht hergeleitet oder numerisch geprüft ist.
	

	
	\section{Einführung}
	\label{sec:introduction}
	
	Die Teilchenphysik steht vor einem fundamentalen Problem: Das Standardmodell mit seinen über zwanzig freien Parametern bietet keine Erklärung für die beobachteten Teilchenmassen. Diese erscheinen willkürlich und ohne theoretische Rechtfertigung. Das T0-Modell setzt dem zwei komplementäre Berechnungsmethoden entgegen, die mit dem einen geometrischen Parameter $\xi$ arbeiten und auch einen Ansatz für die Neutrino-Massen enthalten.
	
	\subsection{Das Parameter-Problem des Standardmodells}
	\label{subsec:parameter_problem}
	
	Das Standardmodell hat trotz seines experimentellen Erfolgs eine theoretische Lücke: Es enthält mehr als 20 freie Parameter, die experimentell bestimmt werden müssen. Diese umfassen:
	
	\begin{itemize}
		\item \textbf{Fermion-Massen}: 9 geladene Lepton- und Quark-Massen
		\item \textbf{Neutrino-Massen}: 3 Neutrino-Masseneigenwerte
		\item \textbf{Mischungsparameter}: 4 CKM- und 4 PMNS-Matrix-Elemente
		\item \textbf{Eichkopplungen}: 3 fundamentale Kopplungskonstanten
		\item \textbf{Higgs-Parameter}: Vakuumerwartungswert und Selbstkopplung
		\item \textbf{QCD-Parameter}: Starke CP-Phase und andere
	\end{itemize}
	
	\begin{important}[Parameterreduktion]
		Das T0-Modell führt die über zwanzig freien Parameter des Standardmodells auf den einen geometrischen Parameter $\xipar = \frac{4}{3} \times 10^{-4}$ zurück, der aus der Geometrie des dreidimensionalen Raums folgt. Dazu kommen rationale Koeffizienten $r_i$, die in der Etablierungsphase an den gemessenen Massen festgelegt wurden (siehe den Abschnitt zur methodischen Klarstellung). Diese Version enthält zusätzlich einen Ansatz für die Neutrino-Quantenzahlen sowie die quantenfeldtheoretische Herleitung.
	\end{important}
	
	\section{Methodische Klarstellung: Etablierung vs. Vorhersage}
	\label{sec:methodische_klarstellung}
	
	\begin{important}[Wissenschaftshistorische Einordnung]
		Das T0-Modell orientiert sich an der Methodik der \textbf{Muster-Erkennung und systematischen Klassifikation}, analog zur Entwicklung des Periodensystems (Mendeleev 1869) oder des Quark-Modells (Gell-Mann 1964).
	\end{important}
	
	\subsection{Zwei-Phasen-Entwicklung}
	\label{subsec:zwei_phasen}
	
	\textbf{Phase 1: Etablierung der Systematik}
	\begin{enumerate}
		\item Muster-Erkennung in bekannten Teilchenmassen (Elektron, Myon, Tau)
		\item Parameter-Bestimmung aus experimentellen Daten
		\item Quantenzahl-Zuordnung etablieren
		\item Mathematische Äquivalenz beider Methoden zeigen
	\end{enumerate}
	
	\textbf{Phase 2: Vorhersagekraft entfalten}
	\begin{enumerate}
		\item Extrapolation auf unbekannte Teilchen
		\item Quark-Sektor aus Lepton-Mustern ableiten
		\item Neue Generationen vorhersagen
		\item Experimentelle Tests durchführen
	\end{enumerate}
	
	\subsection{Historische Präzedenz erfolgreicher Muster-Physik}
	\label{subsec:historische_praezedenz}
	
	Das T0-Modell orientiert sich an der Methodik früherer Muster-Entdeckungen:
	
	\begin{table}[H]
		\centering
		\adjustbox{max width=\textwidth}{%
\begin{tabular}{p{3cm}p{4cm}p{4cm}p{3cm}}
			\toprule
			\textbf{Entdeckung} & \textbf{Muster-Erkennung} & \textbf{Vorhersagen} & \textbf{Bestätigung} \\
			\midrule
			Periodensystem (1869) & Atomgewichte und Eigenschaften & Gallium, Germanium, Scandium & Experimentell bestätigt \\
			Spektrallinien (1885) & Wasserstoff-Linien & Rydberg-Formel für alle Serien & Quantenmechanik \\
			Quark-Modell (1964) & Hadron-Massen & Achtfacher Weg & QCD-Theorie \\
			\textbf{T0-Modell (2025)} & \textbf{Lepton-Massen} & \textbf{4. Generation, Quarks} & \textbf{Experimentelle Tests} \\
			\bottomrule
		\end{tabular}%
}
		\caption{Historische Präzedenz der Muster-Physik}
		\label{tab:historische_praezedenz}
	\end{table}
	
	\section{Von Energiefeldern zu Teilchenmassen}
	\label{sec:energy_fields_to_masses}
	
	\subsection{Die fundamentale Herausforderung}
	\label{subsec:fundamental_challenge}
	
	Ein Kernanliegen des T0-Modells ist, Teilchenmassen aus geometrischen Prinzipien zu berechnen. Während das Standardmodell über 20 freie Parameter zur Beschreibung von Teilchenmassen benötigt, arbeitet das T0-Modell mit der geometrischen Konstante $\xigeom = \frac{4}{3} \times 10^{-4}$ und rationalen Koeffizienten $r_i$.
	
	\begin{tcolorbox}[colback=green!5!white,colframe=green!75!black,title=Massenbeschreibung]
		\textbf{Parameter-Reduktion:}
		\begin{itemize}
			\item \textbf{Standardmodell}: 20+ freie Massenparameter (empirisch bestimmt)
			\item \textbf{T0-Modell}: ein geometrischer Parameter $\xi$, dazu rationale Koeffizienten $r_i$
			\item \textbf{Experimentelle Genauigkeit}: 99\% durchschnittliche Übereinstimmung (geladene Fermionen; Neutrinos nur mit Obergrenzen verglichen)
			\item \textbf{Theoretische Grundlage}: Dreidimensionale Raumgeometrie + QFT-Herleitung
		\end{itemize}
	\end{tcolorbox}
	
	\subsection{Energiebasiertes Massenkonzept}
	\label{subsec:energy_based_mass}
	
	Im T0-Framework wird angenommen, dass das, was wir traditionell als ``Masse'' bezeichnen, eine Manifestation charakteristischer Energieskalen von Feldanregungen ist:
	
	\begin{equation}
		\boxed{m_i \rightarrow E_{\text{char},i} \quad \text{(charakteristische Energie von Teilchentyp } i\text{)}}
		\label{eq:mass_to_energy}
	\end{equation}
	
	Diese Umdeutung hebt die Unterscheidung zwischen Masse und Energie auf und behandelt beide als verschiedene Aspekte derselben fundamentalen Größe.
	
	\section{Zwei komplementäre Berechnungsmethoden}
	\label{sec:two_calculation_methods}
	
	Das T0-Modell bietet zwei mathematisch äquivalente, aber konzeptionell verschiedene Ansätze zur Berechnung von Teilchenmassen:
	
	\subsection{Methode 1: Direkte geometrische Resonanz}
	\label{subsec:direct_geometric_method}
	
	\textbf{Konzeptionelle Grundlage:} Teilchen als Resonanzen im universellen Energiefeld
	
	Die direkte Methode behandelt Teilchen als charakteristische Resonanzmoden des Energiefelds $\Efield$, analog zu stehenden Wellenmustern:
	
	\begin{equation}
		\text{Teilchen} = \text{Diskrete Resonanzmoden von } \Efield(x,t)
	\end{equation}
	
	\textbf{Drei-Schritt-Berechnungsprozess:}
	
	\textbf{Schritt 1: Geometrische Quantisierung}
	\begin{equation}
		\xi_i = \xi_0 \cdot f(n_i, l_i, j_i)
		\label{eq:geometric_quantization}
	\end{equation}
	
	wobei:
	\begin{align}
		\xi_0 &= \frac{4}{3} \times 10^{-4} \quad \text{(geometrischer Basisparameter)} \\
		n_i, l_i, j_i &= \text{Quantenzahlen aus 3D-Wellengleichung} \\
		f(n_i, l_i, j_i) &= \text{geometrische Funktion aus räumlichen Harmonien}
	\end{align}
	
	\textbf{Schritt 2: Resonanzfrequenzen}
	\begin{equation}
		\omega_i = \frac{c^2}{\xi_i \cdot r_{\text{char}}}
		\label{eq:resonance_frequencies}
	\end{equation}
	
	In natürlichen Einheiten ($c = 1$):
	\begin{equation}
		\omega_i = \frac{1}{\xi_i}
	\end{equation}
	
	\textbf{Schritt 3: Massenbestimmung aus Energieerhaltung}
	\begin{equation}
		E_{\text{char},i} = \hbar \omega_i = \frac{\hbar}{\xi_i}
		\label{eq:energy_from_frequency}
	\end{equation}
	
	In natürlichen Einheiten ($\hbar = 1$):
	\begin{equation}
		\boxed{E_{\text{char},i} = \frac{1}{\xi_i}}
		\label{eq:characteristic_energy_direct}
	\end{equation}
	
	\subsection{Methode 2: Erweiterte Yukawa-Methode}
	\label{subsec:extended_yukawa_method}
	
	\textbf{Konzeptionelle Grundlage:} Brücke zur Standardmodell-Formulierung
	
	Die erweiterte Yukawa-Methode behält die Kompatibilität mit Standardmodell-Berechnungen bei, während sie die Yukawa-Kopplungen auf eine geometrische Struktur zurückführt:
	
	\begin{equation}
		E_{\text{char},i} = y_i \cdot v
		\label{eq:yukawa_mass_formula}
	\end{equation}
	
	wobei $v = 246$ GeV der Higgs-Vakuumerwartungswert ist.
	
	\textbf{Geometrische Yukawa-Kopplungen:}
	\begin{equation}
		\boxed{y_i = r_i \cdot \left(\frac{4}{3} \times 10^{-4}\right)^{\pi_i}}
		\label{eq:geometric_yukawa}
	\end{equation}
	
	\textbf{Generationshierarchie:}
	\begin{align}
		\text{1. Generation:} \quad &\pi_i = \frac{3}{2} \quad \text{(Elektron, Up-Quark)} \\
		\text{2. Generation:} \quad &\pi_i = 1 \quad \text{(Myon, Charm-Quark)} \\
		\text{3. Generation:} \quad &\pi_i = \frac{2}{3} \quad \text{(Tau, Top-Quark)}
	\end{align}
	
	Die Koeffizienten $r_i$ sind einfache rationale Zahlen, die der geometrischen Struktur jedes Teilchentyps zugeordnet werden.
	
	\section{Quantenfeldtheoretische Herleitung der $\xi$-Konstante}
	\label{sec:qft_herleitung}
	
	\subsection{EFT-Matching und Yukawa-Kopplung nach EWSB}
	\label{subsec:eft_matching}
	
	Nach der elektroschwachen Symmetriebrechung haben wir die Yukawa-Wechselwirkung:
	
	\begin{equation}
		\mathcal{L}_{\text{Yukawa}} \supset -\lambda_h \bar{\psi}\psi H, \quad \text{mit} \quad H = \frac{v + h}{\sqrt{2}}
	\end{equation}
	
	Nach EWSB:
	\begin{equation}
		\mathcal{L} \supset -m \bar{\psi}\psi - y h \bar{\psi}\psi
	\end{equation}
	
	mit den Beziehungen:
	\begin{equation}
		m = \frac{\lambda_h v}{\sqrt{2}} \quad \text{und} \quad y = \frac{\lambda_h}{\sqrt{2}}
	\end{equation}
	
	Die lokale Massenabhängigkeit auf das physikalische Higgs-Feld $h(x)$ führt zu:
	
	\begin{equation}
		m(h) = m\left(1 + \frac{h}{v}\right) \quad \Rightarrow \quad \partial_\mu m = \frac{m}{v}\partial_\mu h
	\end{equation}
	
	\subsection{T0-Operatoren in der effektiven Feldtheorie}
	\label{subsec:t0_operators}
	
	In der T0-Theorie treten Operatoren der Form auf:
	
	\begin{equation}
		O_T = \bar{\psi}\gamma^\mu\Gamma_\mu^{(T)}\psi
	\end{equation}
	
	mit dem charakteristischen Zeitfeld-Kopplungsterm:
	\begin{equation}
		\Gamma_\mu^{(T)} = \frac{\partial_\mu m}{m^2}
	\end{equation}
	
	Einsetzen der Higgs-Abhängigkeit:
	\begin{equation}
		\Gamma_\mu^{(T)} = \frac{\partial_\mu m}{m^2} = \frac{1}{mv}\partial_\mu h
	\end{equation}
	
	Dies legt nahe, dass ein $\partial_\mu h$-gekoppelter Vektorstrom der UV-Ursprung ist.
	
	\subsection{1-Loop-Matching-Rechnung}
	\label{subsec:one_loop_matching}
	
	Die 1-Loop-Amplitude für den T0-Vertex ergibt:
	\begin{equation}
		F_V(0) = \frac{y^2}{16\pi^2}\left[\frac{1}{2} - \frac{1}{2}\ln\left(\frac{m_h^2}{\mu^2}\right) + \frac{r(r-\ln r-1)}{(r-1)^2}\right]
	\end{equation}
	
	Für hierarchische Massen ($m \ll m_h$) dominiert der konstante Term:
	\begin{equation}
		F_V(0) \approx \frac{y^2}{32\pi^2}
	\end{equation}
	
	\subsection{Finale $\xi$-Formel aus Higgs-Physik}
	\label{subsec:finale_xi_formel}
	
	Das EFT-Matching liefert die fundamentale Beziehung:
	\begin{equation}
		\boxed{\xi = \frac{\lambda_h^2 v^2}{16\pi^3 m_h^2}}
	\end{equation}
	
	Mit Standard-Higgs-Parametern ($m_h = 125.1$ GeV, $v = 246.22$ GeV, $\lambda_h \approx 0.13$):
	\begin{equation}
		\xi \approx 1.318 \times 10^{-4}
	\end{equation}
	
	Dies stimmt gut mit der geometrischen Bestimmung $\xi_0 = \frac{4}{3} \times 10^{-4} \approx 1.333 \times 10^{-4}$ überein (Abweichung $\approx 1.15\%$).
	
	\section{Universelle Teilchenmassen-Systematik}
	\label{sec:universal_masses}
	
	\subsection{Überarbeitete Universaltabelle der Fermionen}
	\label{subsec:universal_table}
	
	\adjustbox{max width=\textwidth}{%
\begin{tabular}{|l|c|c|c|c|c|l|}
		\hline
		Fermion & Generation & Family & Spin & $r_f$ & Exponent $p_f$ & Symmetrie \\
		\hline
		Electron Neutrino & 1 & 0 & 1/2 & $4/3$ & $5/2$ & Doppeltes $\xi$ \\
		Electron          & 1 & 0 & 1/2 & $4/3$  & $3/2$ & Leptonenzahl \\
		Muon Neutrino     & 2 & 1 & 1/2 & $16/5$ & $2$   & Doppeltes $\xi$ \\
		Muon              & 2 & 1 & 1/2 & $16/5$ & $1$   & Leptonenzahl \\
		Tau Neutrino      & 3 & 2 & 1/2 & $25/9$ & $5/3$ & Doppeltes $\xi$ \\
		Tau               & 3 & 2 & 1/2 & $25/9$ & $2/3$ & Leptonenzahl \\
		\hline
		Up     & 1 & 0 & 1/2 & $6$          & $3/2$ & Color \\
		Down   & 1 & 0 & 1/2 & $\tfrac{25}{2}$ & $3/2$ & Color + Isospin \\
		Charm  & 2 & 1 & 1/2 & $2$$^*$          & $2/3$ & Color \\
		Strange& 2 & 1 & 1/2 & $\tfrac{26}{9}$ & $1$   & Color \\
		Top    & 3 & 2 & 1/2 & $\tfrac{1}{28}$ & $-1/3$ & Color \\
		Bottom & 3 & 2 & 1/2 & $\tfrac{3}{2}$  & $1/2$ & Color \\
		\hline
	\end{tabular}%
}
	
	\footnotetext{* Korrigiert von ursprünglich $8/9$ basierend auf detaillierter numerischer Analyse}
	
	\textit{(Korrigiert am 30.~Sept.~2026: vorher Neutrino-Exponenten $3$ ($\nu_\mu$) und $25/9$ ($\nu_\tau$) -- nach Gl.~\eqref{eq:neutrino_suppression} ($f_\nu = f_{\text{Lepton}}\cdot\xi$) gilt $p_\nu = p_\ell + 1$, also $5/2$, $2$, $5/3$; $25/9$ ist der Koeffizient $r_\tau$.)}
	
	\section{Vollständige numerische Rekonstruktion}
	\label{sec:vollstaendige_rekonstruktion}
	
	Die folgende Analyse zeigt die explizite Berechnung aller Fermionen mit beiden Methoden:
	
	\subsection{Grundlagen und experimentelle Eingangsdaten}
	\label{subsec:grundlagen}
	
	\textbf{Fundamentale Konstanten:}
	\begin{align}
		\xi_0 = \xi &= \frac{4}{3} \times 10^{-4} = 1.333333333\ldots \times 10^{-4} \\
		v &= 246 \text{ GeV}
	\end{align}
	
	\textbf{Experimentelle Massen (PDG-nahe Werte):}
	\begin{align}
		m_e^{\text{exp}} &= 0.0005109989461 \text{ GeV} \\
		m_\mu^{\text{exp}} &= 0.1056583745 \text{ GeV} \\
		m_\tau^{\text{exp}} &= 1.77686 \text{ GeV}
	\end{align}
	
	\subsection{Geladene Leptonen: Detaillierte Berechnungen}
	\label{subsec:charged_leptons_detailed}
	
	\textbf{Elektronmassen-Berechnung:}
	
	\textit{Direkte Methode:}
	\begin{align}
		\xi_e &= \frac{4}{3} \times 10^{-4} \times f_e(1,0,1/2) \\
		&= \frac{4}{3} \times 10^{-4} \times 1 = \frac{4}{3} \times 10^{-4} \\
		E_{e} &= \frac{1}{\xi_e} = \frac{3}{4 \times 10^{-4}} = 0.511 \text{ MeV}
	\end{align}
	\textit{(Vermerk vom 30.~Sept.~2026: $1/\xi_e = 3/(4\times10^{-4}) = 7500$ ist eine reine Zahl; der Schritt zu $0.511$~MeV ist nicht angegeben und einheitenabhängig. Derselbe Fehler ist für Dok.~006 im Korrekturregister (R109) beschrieben; maßgeblich ist die Yukawa-Darstellung $m_i = r_i\,\xi^{p_i}\,v$.)}
	
	\textit{Erweiterte Yukawa-Methode:}
	\begin{align}
		r_e &= \frac{m_e^{\text{exp}}}{v \cdot \xi^{3/2}} \approx 1.349 \\
		y_e &= 1.349 \times \left(\frac{4}{3} \times 10^{-4}\right)^{3/2} \\
		E_e &= y_e \times 246 \text{ GeV} = 0.511 \text{ MeV}
	\end{align}
	
	\textbf{Myonmassen-Berechnung:}
	
	\textit{Direkte Methode:}
	\begin{align}
		\xi_\mu &= \frac{4}{3} \times 10^{-4} \times f_\mu(2,1,1/2) \\
		&= \frac{4}{3} \times 10^{-4} \times \frac{16}{5} = \frac{64}{15} \times 10^{-4} \\
		E_{\mu} &= \frac{1}{\xi_\mu} = 105.66 \text{ MeV}
	\end{align}
	\textit{(Vermerk vom 30.~Sept.~2026: $1/\xi_\mu = 15/(64\times10^{-4}) = 2343.75$ ist eine reine Zahl; nach dieser Formel wäre das Myon um den Faktor $16/5$ leichter als das Elektron ($7500/2343.75 = 3.2$). Der Schritt zu $105.66$~MeV ist nicht angegeben; vgl. R109 zu Dok.~006.)}
	
	\textit{Erweiterte Yukawa-Methode:}
	\begin{align}
		y_\mu &= \frac{16}{5} \times \left(\frac{4}{3} \times 10^{-4}\right)^1 = 4.267 \times 10^{-4} \\
		E_\mu &= y_\mu \times 246 \text{ GeV} = 104.96 \text{ MeV}
	\end{align}
	\textbf{Experiment:} $105.66 \text{ MeV}$ → Abweichung $\approx 0.65\%$
	
	\subsection{Vollständige Neutrino-Behandlung}
	\label{sec:complete_neutrino_treatment}
	
	\begin{neutrino}[Neutrino-Ansatz]
		\textbf{[S]} Das T0-Modell enthält einen geometrischen Ansatz für die Neutrino-Massen über die Annahme einer charakteristischen \textbf{doppelten $\xi$-Unterdrückung}. Damit wird die bisherige Lücke im Neutrino-Sektor adressiert; der Vergleich mit den Oszillationsdaten steht noch aus.
	\end{neutrino}
	
	\subsection{Neutrino-Quantenzahlen}
	\label{subsec:neutrino_quantum_numbers}
	
	Neutrinos folgen derselben Quantenzahl-Struktur wie andere Fermionen, aber mit einer wesentlichen Modifikation aufgrund ihrer schwachen Wechselwirkungsnatur:
	
	\begin{table}[H]
		\centering
		\adjustbox{max width=\textwidth}{%
\begin{tabular}{lcccc}
			\toprule
			\textbf{Neutrino} & \textbf{n} & \textbf{l} & \textbf{j} & \textbf{Unterdrückung} \\
			\midrule
			$\nu_e$ & 1 & 0 & 1/2 & Doppeltes $\xi$ \\
			$\nu_\mu$ & 2 & 1 & 1/2 & Doppeltes $\xi$ \\
			$\nu_\tau$ & 3 & 2 & 1/2 & Doppeltes $\xi$ \\
			\bottomrule
		\end{tabular}%
}
		\caption{Neutrino-Quantenzahlen mit charakteristischer doppelter $\xi$-Unterdrückung}
		\label{tab:neutrino_quantum_numbers}
	\end{table}
	
	\subsection{Doppelte $\xi$-Unterdrückungsmechanismus}
	\label{subsec:double_xi_suppression}
	
	Die zentrale Annahme ist, dass Neutrinos einen zusätzlichen geometrischen Unterdrückungsfaktor erfahren:
	
	\begin{equation}
		f(n_{\nu_i}, l_{\nu_i}, j_{\nu_i}) = f(n_i, l_i, j_i)_{\text{Lepton}} \times \xi
		\label{eq:neutrino_suppression}
	\end{equation}
	
	\textbf{Neutrino-Massenberechnungen:}
	
	\textbf{Elektron-Neutrino:}
	\begin{align}
		\xi_{\nu_e} &= \frac{4}{3} \times 10^{-4} \times 1 \times \frac{4}{3} \times 10^{-4} = \frac{16}{9} \times 10^{-8} \\
		E_{\nu_e} &= \frac{1}{\xi_{\nu_e}} = 9.1 \text{ meV}
	\end{align}
	
	\textbf{Myon-Neutrino:}
	\begin{align}
		\xi_{\nu_\mu} &= \frac{4}{3} \times 10^{-4} \times \frac{16}{5} \times \frac{4}{3} \times 10^{-4} = \frac{256}{45} \times 10^{-8} \\
		E_{\nu_\mu} &= \frac{1}{\xi_{\nu_\mu}} = 1.9 \text{ meV}
	\end{align}
	
	\textbf{Tau-Neutrino:}
	\begin{align}
		\xi_{\nu_\tau} &= \frac{4}{3} \times 10^{-4} \times \frac{25}{9} \times \frac{4}{3} \times 10^{-4} = \frac{400}{81} \times 10^{-8} \\
		E_{\nu_\tau} &= \frac{1}{\xi_{\nu_\tau}} = 18.0 \text{ meV}
	\end{align}
	
	\textit{(Vermerk vom 30.~Sept.~2026: Die Werte $9.1$, $1.9$ und $18.0$~meV sind dieselben wie im Neutrinoteil von Dok.~006, der im Korrekturregister (R109) als nicht tragfähig geführt ist. Für den Schritt von $1/\xi_\nu$ zu meV gibt es keinen gemeinsamen Anker, und die Werte sind mit den Oszillationsdaten unverträglich, die für das schwerste Neutrino mindestens etwa $50$~meV verlangen. Maßgeblich ist Dok.~340.)}
	
	\section{Vollständige Quark-Analyse mit beiden Methoden}
	\label{sec:quark_analyse}
	
	\subsection{Explizite Berechnungen der Quarkmassen}
	\label{subsec:quark_calculations}
	
	Wir verwenden $\xi=\tfrac{4}{3}\times10^{-4}$ und $v=246\ \mathrm{GeV}$.
	Für die Yukawa-Darstellung:
	\[
	y_i = r_i\,\xi^{p_i},\qquad m_i^{\text{pred}}=y_i\,v.
	\]
	Für die direkte geometrische Darstellung:
	\[
	f_i=\frac{1}{\xi\, m_i^{\text{exp}}},\qquad m_i^{\text{exp}}=\frac{1}{\xi\, f_i}.
	\]
	
	\begin{table}[H]
		\centering
		\adjustbox{max width=\textwidth}{%
\begin{tabular}{lcccccc}
			\toprule
			Quark & $p_i$ & $r_i$ (korr.) & $m_i^{\text{pred}}$ & $m_i^{\text{exp}}$ & rel.\ Fehler & Bemerkung\\
			& & & (GeV) & (GeV) & (\%) & \\
			\midrule
			Up     & $3/2$ & $6$        & $2.272\times10^{-3}$ & $2.16\times10^{-3}$ & $+5.21$ & OK \\
			Down   & $3/2$ & $25/2$     & $4.734\times10^{-3}$ & $4.72\times10^{-3}$ & $+0.30$ & OK \\
			Strange& $1$   & $26/9$        & $9.476\times10^{-2}$  & $9.35\times10^{-2}$  & $+1.34$ & OK\\
			Charm  & $2/3$ & $2$      & $1.284\times10^{0}$  & $1.28$              & $+0.32$ & Korrigiert\\
			Bottom & $1/2$ & $3/2$      & $4.261\times10^{0}$   & $4.183$              & $+1.86$ & OK \\
			Top    & $-1/3$& $1/28$     & $1.7197\times10^{2}$  & $172.57$               & $-0.35$ & OK \\
			\bottomrule
		\end{tabular}%
}
		\caption{Yukawa-Vorhersagen mit korrigierten $r_i,p_i$ und Vergleich mit Referenzmassen.}
		\label{tab:quark_yukawa_predictions}
	\end{table}
	
	\textit{(Korrigiert am 30.~Sept.~2026: vorher Strange $9.50\times10^{-2}$~GeV, $0.00\,\%$, \enquote{Exakt} und Charm $1.279$~GeV, $-0.08\,\%$ -- $(26/9)\,\xi\,v = 94.76$~MeV und $2\,\xi^{2/3}\,v = 1.284$~GeV. Die Referenzmassen sind auf PDG~2024 gesetzt (vorher $m_u = 2.27$~MeV, $m_s = 95.0$~MeV, $m_b = 4.26$~GeV, $m_t = 171$~GeV, die mit den Vorhersagen zusammenfielen); die relativen Fehler sind damit neu berechnet.)}
	
	\subsection{Korrektur für das Charm-Quark}
	\label{subsec:charm_correction}
	
	Die ursprünglich in der Tabelle angegebene Größe $r_c=8/9$ reproduziert nicht die referenzierte Masse $m_c=1.28\ \mathrm{GeV}$. Der notwendige Wert ist:
	\[
	r_c^{\text{required}}=\frac{m_c^{\text{exp}}}{v\,\xi^{2/3}}\approx 1.994 \approx 2.
	\]
	
	Daher wurde in der korrigierten Universaltabelle $r_c \approx 2$ eingesetzt.
	
	\section{Umfassende experimentelle Validierung}
	\label{sec:comprehensive_validation}
	
	\subsection{Vollständige Genauigkeitsanalyse}
	\label{subsec:complete_accuracy}
	
	Die folgende Tabelle vergleicht die T0-Werte mit den experimentellen Werten:
	
	\begin{table}[H]
		\centering
		\adjustbox{max width=\textwidth}{%
\begin{tabular}{lcccc}
			\toprule
			\textbf{Teilchen} & \textbf{T0-Vorhersage} & \textbf{Experiment} & \textbf{Genauigkeit} & \textbf{Typ} \\
			\midrule
			\multicolumn{5}{c}{\textit{Geladene Leptonen}} \\
			\midrule
			Elektron & 0.505 MeV & 0.511 MeV & 98.82\% & Lepton \\
			Myon & 104.96 MeV & 105.66 MeV & 99.35\% & Lepton \\
			Tau & 1783.4 MeV & 1776.86 MeV & 99.63\% & Lepton \\
			\midrule
			\multicolumn{5}{c}{\textit{Neutrinos}} \\
			\midrule
			$\nu_e$ & 9.1 meV & $< 450$ meV & Kompatibel & Neutrino \\
			$\nu_\mu$ & 1.9 meV & $< 180$ keV & Kompatibel & Neutrino \\
			$\nu_\tau$ & 18.0 meV & $< 18$ MeV & Kompatibel & Neutrino \\
			\midrule
			\multicolumn{5}{c}{\textit{Quarks}} \\
			\midrule
			Up-Quark & 2.272 MeV & 2.16 MeV & 94.79\% & Quark \\
			Down-Quark & 4.734 MeV & 4.72 MeV & 99.70\% & Quark \\
			Strange-Quark & 94.76 MeV & 93.5 MeV & 98.66\% & Quark \\
			Charm-Quark & 1.284 GeV & 1.28 GeV & 99.68\% & Quark \\
			Bottom-Quark & 4.261 GeV & 4.183 GeV & 98.14\% & Quark \\
			Top-Quark & 171.97 GeV & 172.57 GeV & 99.65\% & Quark \\
			\midrule
			\textbf{Durchschnitt} & & & \textbf{98.7\%} & \textbf{Alle Fermionen} \\
			\bottomrule
		\end{tabular}%
}
		\caption{Vollständige experimentelle Validierung der T0-Modell-Vorhersagen}
		\label{tab:complete_validation}
	\end{table}
	
	\textit{(Korrigiert am 30.~Sept.~2026: vorher Elektron $0.511$~MeV/$99.98\,\%$, Tau $1777.1$~MeV/$99.99\,\%$, Strange $95.0$~MeV/$100.0\,\%$, Charm $1.279$~GeV/$99.92\,\%$, Top $171.99$~GeV, Durchschnitt $99.6\,\%$ -- mit den tabellierten Koeffizienten gilt $\tfrac{4}{3}\,\xi^{3/2}v = 0.505$~MeV ($-1.2\,\%$; $99.98\,\%$ nur mit dem an $m_e$ angepassten $r_e = 1.349$), $\tfrac{25}{9}\,\xi^{2/3}v = 1783.4$~MeV ($+0.37\,\%$), $\tfrac{26}{9}\,\xi v = 94.76$~MeV, $2\,\xi^{2/3}v = 1.284$~GeV, $\tfrac{1}{28}\,\xi^{-1/3}v = 171.97$~GeV; Quark-Referenzmassen nach PDG~2024; Mittelwert der neun geladenen Fermionen $98.7\,\%$.)}
	
	\begin{keyresult}[Zusammenfassung der Genauigkeit]
		Das T0-Modell erreicht 98.7\% durchschnittliche Genauigkeit \textit{(Korrigiert am 30.~Sept.~2026: vorher 99.6\% -- siehe Vermerk unter Tabelle~\ref{tab:complete_validation}.)} über die \textbf{geladenen} Fermionen mit dem einen geometrischen Parameter $\xi$ und den tabellierten rationalen Koeffizienten $r_i$. Der Neutrino-Sektor ist als Ansatz enthalten \textbf{[S]}; er wird bisher nur mit Obergrenzen verglichen.
	\end{keyresult}
	
	\section{Vorhersagekraft des etablierten Systems}
	\label{sec:vorhersagekraft}
	
	\subsection{Neue Teilchen-Generationen}
	\label{subsec:neue_generationen}
	
	Mit den etablierten Mustern können neue Teilchen vorhergesagt werden:
	
	\textbf{4. Generation (extrapoliert) [S]:}
	\begin{align}
		n &= 4, \quad \pi_4 = \frac{1}{2}, \quad r_4 \approx 2.0 \\
		m_{\text{4.Gen}} &= r_4 \times \xi^{1/2} \times v \approx 5.7 \text{ GeV}
	\end{align}
	
	\textit{(Vermerk vom 30.~Sept.~2026: Die Rechnung $2.0 \times \xi^{1/2} \times 246$~GeV $\approx 5.7$~GeV stimmt; ein geladenes Lepton der 4.~Generation bei dieser Masse ist jedoch experimentell ausgeschlossen (LEP, untere Grenze etwa $100$~GeV).)}
	
	\subsection{Quark-Sektor Extrapolation}
	\label{subsec:quark_extrapolation}
	
	Die Lepton-Muster lassen sich auf Quarks übertragen:
	
	\begin{table}[H]
		\centering
		\adjustbox{max width=\textwidth}{%
\begin{tabular}{lcccc}
			\toprule
			\textbf{Quark} & \textbf{Generation} & \textbf{$r_i$} & \textbf{$\pi_i$} & \textbf{Vorhersage} \\
			\midrule
			Up & 1 & 6 & 3/2 & 2.3 MeV \\
			Down & 1 & 12.5 & 3/2 & 4.7 MeV \\
			Charm & 2 & 2.0 & 2/3 & 1.3 GeV \\
			Strange & 2 & 2.89 & 1 & 95 MeV \\
			Top & 3 & 0.036 & -1/3 & 173 GeV \\
			Bottom & 3 & 1.5 & 1/2 & 4.3 GeV \\
			\bottomrule
		\end{tabular}%
}
		\caption{Quark-Vorhersagen aus etablierten Mustern}
		\label{tab:quark_vorhersagen}
	\end{table}
	
	\section{Korrigierte Interpretation der mathematischen Äquivalenz}
	\label{sec:korrigierte_interpretation}
	
	\begin{keypoint}[Bedeutung der Äquivalenz]
		Die mathematische Äquivalenz beider Methoden ist \textbf{per Definition gegeben}, wenn die Parameter ($r_i$ oder $f_i$) aus denselben experimentellen Massen bestimmt werden. Die Äquivalenz ist kein Beweis für die Theorie, sondern eine Konsistenz-Eigenschaft der mathematischen Struktur.
	\end{keypoint}
	
	\subsection{Transformationsbeziehung als Brücke}
	\label{subsec:transformationsbeziehung}
	
	Die fundamentale Beziehung:
	\begin{equation}
		f_i = \frac{1}{r_i \, \xi^{\pi_i} \, v \, \xi_0}
		\label{eq:transformation_bridge}
	\end{equation}
	
	verknüpft beide Methoden mathematisch. Wenn $r_i$ aus experimentellen Massen bestimmt wird, folgt $f_i$ automatisch und umgekehrt.
	
	\begin{table}[H]
		\centering
		\adjustbox{max width=\textwidth}{%
\begin{tabular}{lcccc}
			\toprule
			\textbf{Teilchen} & \textbf{$m^{\text{exp}}$ (GeV)} & \textbf{$r_i$ (Yukawa)} & \textbf{$f_i$ (direkt)} & \textbf{Genauigkeit} \\
			\midrule
			Elektron & 0.000511 & 1.349 & $1.468 \times 10^{7}$ & $99.98\%$ \\
			Myon & 0.10566 & 3.221 & $7.099 \times 10^{4}$ & $99.35\%$ \\
			Tau & 1.77686 & $25/9 = 2.778$ & $4.221 \times 10^{3}$ & $99.63\%$ \\
			\midrule
			$\nu_e$ & 9.1 $\times 10^{-12}$ & 1.349 & $8.24 \times 10^{14}$ & Vorhersage \\
			$\nu_\mu$ & 1.9 $\times 10^{-12}$ & 3.221 & $3.95 \times 10^{15}$ & Vorhersage \\
			$\nu_\tau$ & 18.0 $\times 10^{-12}$ & 2.778 & $4.17 \times 10^{14}$ & Vorhersage \\
			\bottomrule
		\end{tabular}%
}
		\caption{Numerische Äquivalenz beider T0-Methoden für alle Leptonen}
		\label{tab:numerische_aequivalenz_komplett}
	\end{table}
	
	\textit{(Korrigiert am 30.~Sept.~2026: Tau vorher $r_i = 2.778$ mit $99.99\,\%$ -- die Genauigkeit gehörte zum aus dem Messwert gefitteten $r_\tau = 2.7675$; mit dem rationalen $r_\tau = 25/9$ ergibt die Leiter $1.7834$~GeV, also $99.63\,\%$. Elektron und Myon stehen hier mit den aus dem Messwert bestimmten $r_i$; die rationalen Werte sind $4/3$ und $16/5$.)}
	
	\textit{(Korrigiert am 30.~Sept.~2026: vorher $m^{\text{exp}} = 9.1/1.9/18.0 \times 10^{-6}$~GeV und $f_i = 8.235\times10^{10}$, $3.947\times10^{11}$, $3.989\times10^{10}$ -- $1$~meV $= 10^{-12}$~GeV; mit $f_i = 1/(\xi\,m_i)$ folgt $8.24\times10^{14}$, $3.95\times10^{15}$, $4.17\times10^{14}$. Die alten $f_i$ folgen auch aus den alten Massen nicht ($8.24\times10^{8}$, $3.95\times10^{9}$, $4.17\times10^{8}$). Zu den meV-Werten selbst siehe den Vermerk im Neutrino-Abschnitt.)}
	
	\section{Experimentelle Vorhersagen und Präzisionstests}
	\label{sec:experimentelle_vorhersagen}
	
	\subsection{Modifizierte QED-Vertex-Korrekturen}
	\label{subsec:qed_corrections}
	
	Die T0-Theorie sagt modifizierte Feynman-Regeln voraus:
	\begin{align}
		\text{Zeitfeld-Vertex:} \quad &-i\gamma^\mu\Gamma_\mu^{(T)} = i\gamma^\mu\frac{\partial_\mu m}{m^2} \\
		\text{Modifizierter Fermion-Propagator:} \quad &S_F^{(T0)}(p) = S_F(p) \cdot \left[1 + \frac{\beta}{p^2}\right]
	\end{align}
	
	\subsection{Neutrino-Validierung}
	\label{subsec:neutrino_validation}
	
	Die T0-Neutrino-Vorhersagen erfüllen die in der Tabelle genannten Grenzen; ein Vergleich mit den gemessenen Massendifferenzen der Oszillationsexperimente steht in diesem Dokument aus \textbf{[S]}:
	
	\begin{table}[H]
		\centering
		\adjustbox{max width=\textwidth}{%
\begin{tabular}{lccc}
			\toprule
			\textbf{Parameter} & \textbf{T0-Vorhersage} & \textbf{Experimentelle Grenze} & \textbf{Status} \\
			\midrule
			$m_{\nu_e}$ & 9.1 meV & $< 450$ meV (KATRIN) & Erfüllt \\
			$m_{\nu_\mu}$ & 1.9 meV & $< 180$ keV (indirekt) & Erfüllt \\
			$m_{\nu_\tau}$ & 18.0 meV & $< 18$ MeV (indirekt) & Erfüllt \\
			$\sum m_\nu$ & 29.0 meV & $< 60$ meV (Kosmologie 2024) & Erfüllt \\
			\bottomrule
		\end{tabular}%
}
		\caption{T0-Neutrino-Vorhersagen vs. experimentelle Beschränkungen}
		\label{tab:neutrino_validation}
	\end{table}
	
	\textit{(Korrigiert am 30.~Sept.~2026: vorher $\sum m_\nu = 29.8$~meV -- die Summe der drei Tabellenwerte $9.1 + 1.9 + 18.0$ ist $29.0$~meV. Zu den Einzelwerten siehe den Vermerk bei der Neutrino-Massenberechnung.)}
	
	\begin{important}[Neutrino-Massenhierarchie]
		\textbf{[S]} Das T0-Modell ordnet die Massen als $m_{\nu_\mu} < m_{\nu_e} < m_{\nu_\tau}$. Wie diese Zuordnung zu Flavour-Zuständen mit der von Oszillationsdaten bevorzugten \textbf{normalen Ordnung} der Masseneigenzustände zusammenhängt, ist noch zu klären.
	\end{important}
	
	\section{Methodische Fundierung}
	\label{sec:wissenschaftliche_legitimitaet}
	
	\subsection{Umkehrbarkeit des etablierten Systems}
	\label{subsec:umkehrbarkeit}
	
	Nach der Etablierungsphase soll das T0-System vorhersagend werden:
	
	\textbf{Etablierte Lepton-Muster:}
	\begin{align}
		\text{1. Generation (n=1):} \quad &\pi_i = \frac{3}{2}, \quad r_e \approx 1.35 \\
		\text{2. Generation (n=2):} \quad &\pi_i = 1, \quad r_\mu \approx 3.2 \\
		\text{3. Generation (n=3):} \quad &\pi_i = \frac{2}{3}, \quad r_\tau \approx 2.8
	\end{align}
	
	\subsection{Experimentelle Testbarkeit}
	\label{subsec:experimentelle_testbarkeit}
	
	Die T0-Vorhersagen sind experimentell falsifizierbar:
	
	\begin{enumerate}
		\item \textbf{LHC-Suchen:} Neue Teilchen bei charakteristischen Energien (5-6 GeV Bereich)
		\item \textbf{Präzisionsmessungen:} Verfeinerung der $r_i$-Parameter
		\item \textbf{Neutrino-Tests:} Direkte Neutrino-Massenmessungen
	\end{enumerate}
	
	Für das T0-Verfahren sprechen:
	
	\begin{enumerate}
		\item \textbf{Systematische Struktur:} Alle Parameter folgen erkennbaren Mustern
		\item \textbf{Vorhersagekraft:} Nach Etablierung werden neue Teilchen vorhersagbar
		\item \textbf{Experimentelle Testbarkeit:} Vorhersagen sind falsifizierbar
		\item \textbf{QFT-Fundierung:} Quantenfeldtheoretische Herleitung der $\xi$-Konstante
		\item \textbf{Historische Präzedenz:} Methodik der Muster-Physik (Periodensystem, Quark-Modell)
	\end{enumerate}
	
	\section{Parameterstruktur und universelle Systematik}
	\label{sec:parameterfreie_natur}
	
	\begin{important}[Rolle des Parameters $\xi$]
		Der Parameter $\xi$ lässt sich zusätzlich über Higgs-Parameter abschätzen:
		\begin{equation}
			\xi = \frac{\lambda_h^2 v^2}{16\pi^3 m_h^2} \approx 1.318 \times 10^{-4}
		\end{equation}
		Damit ist $\xi$ nicht frei wählbar. Die Koeffizienten $r_i$ sind rationale Zahlen; ihre Herleitung aus der Geometrie ohne Rückgriff auf gemessene Massen ist der nächste Schritt.
	\end{important}
	
	\subsection{Universelle Quantenzahlen-Tabelle}
	\label{subsec:universal_quantum_table}
	
	\begin{table}[H]
		\centering
		\adjustbox{max width=\textwidth}{%
\begin{tabular}{lcccccc}
			\toprule
			\textbf{Teilchen} & \textbf{n} & \textbf{l} & \textbf{j} & \textbf{$r_i$} & \textbf{$p_i$} & \textbf{Speziell} \\
			\midrule
			\multicolumn{7}{c}{\textit{Geladene Leptonen}} \\
			\midrule
			Elektron & 1 & 0 & 1/2 & 4/3 & 3/2 & -- \\
			Myon & 2 & 1 & 1/2 & 16/5 & 1 & -- \\
			Tau & 3 & 2 & 1/2 & 25/9 & 2/3 & -- \\
			\midrule
			\multicolumn{7}{c}{\textit{Neutrinos}} \\
			\midrule
			$\nu_e$ & 1 & 0 & 1/2 & 4/3 & 5/2 & Doppeltes $\xi$ \\
			$\nu_\mu$ & 2 & 1 & 1/2 & 16/5 & 2 & Doppeltes $\xi$ \\
			$\nu_\tau$ & 3 & 2 & 1/2 & 25/9 & 5/3 & Doppeltes $\xi$ \\
			\midrule
			\multicolumn{7}{c}{\textit{Quarks}} \\
			\midrule
			Up & 1 & 0 & 1/2 & 6 & 3/2 & Farbe \\
			Down & 1 & 0 & 1/2 & 25/2 & 3/2 & Farbe + Isospin \\
			Charm & 2 & 1 & 1/2 & 2 & 2/3 & Farbe \\
			Strange & 2 & 1 & 1/2 & 26/9 & 1 & Farbe \\
			Top & 3 & 2 & 1/2 & 1/28 & -1/3 & Farbe \\
			Bottom & 3 & 2 & 1/2 & 3/2 & 1/2 & Farbe \\
			\bottomrule
		\end{tabular}%
}
		\caption{Vollständige universelle Quantenzahlen-Tabelle für alle Fermionen}
		\label{tab:universal_quantum_numbers}
	\end{table}
	
	\textit{(Korrigiert am 30.~Sept.~2026: vorher $p_{\nu_\mu} = 3$ und $p_{\nu_\tau} = 25/9$ -- nach Gl.~\eqref{eq:neutrino_suppression} gilt $p_\nu = p_\ell + 1$, also $5/2$, $2$, $5/3$.)}
	
	\section{Kritische Bewertung und Limitationen}
	\label{sec:kritische_bewertung}
	
	\subsection{Theoretische Offene Fragen}
	\label{subsec:offene_fragen}
	
	\begin{enumerate}
		\item \textbf{Generationsanzahl:} Warum genau drei Generationen plus vierte Vorhersage?
		\item \textbf{Hierarchie-Problem:} Verbindung zwischen verschiedenen Energieskalen
		\item \textbf{CP-Verletzung:} Einbindung der CKM- und PMNS-Mischungsmatrizen
	\end{enumerate}
	
	\section{Abschließende Bewertung}
	\label{sec:abschliessende_bewertung}
	
	\subsection{Wissenschaftlicher Status}
	\label{subsec:wissenschaftlicher_status}
	
	Das T0-Modell bietet einen systematischen Ansatz zur Beschreibung von Teilchenmassen. Die Kombination aus:
	
	\begin{itemize}
		\item \textbf{Hoher numerischer Genauigkeit} (98.7\% über die geladenen Fermionen) \textit{(Korrigiert am 30.~Sept.~2026: vorher 99.6\%, vgl. Tabelle~\ref{tab:complete_validation}.)}
		\item \textbf{Parameterreduktion} (ein geometrischer Parameter $\xi$, rationale Koeffizienten $r_i$)
		\item \textbf{Universeller Abdeckung} (alle bekannten Fermionen)
		\item \textbf{QFT-Konsistenz} (1-Loop-Herleitung der $\xi$-Konstante)
		\item \textbf{Experimenteller Testbarkeit} (spezifische falsifizierbare Vorhersagen)
	\end{itemize}
	
	rechtfertigt eine ernsthafte wissenschaftliche Betrachtung.
	
	\subsection{Bedeutung für die fundamentale Physik}
	\label{subsec:bedeutung_physik}
	
	Falls experimentell bestätigt, würde das T0-Modell unser Verständnis der Teilchenphysik in folgenden Punkten verändern:
	
	\begin{enumerate}
		\item \textbf{Geometrische Interpretation:} Teilchenmassen als Manifestationen der 3D-Raumgeometrie
		\item \textbf{Vereinheitlichung:} Alle Fermionen folgen derselben universellen Struktur
		\item \textbf{Vorhersagekraft:} Neue Teilchen werden aus etablierten Mustern vorhersagbar
		\item \textbf{Theoretische Einfachheit:} Vereinfachung komplexer Phänomene
	\end{enumerate}
	
	Das T0-Modell ist ein Beispiel dafür, dass eine grundlegende Beschreibung nicht unbedingt in größerer Komplexität liegen muss; sie kann auch in einer Vereinfachung liegen.
	
	\begin{thebibliography}{99}
		\bibitem{pascher_t0_energie_2025}
		Pascher, J. (2025). \textit{Das T0-Modell (Planck-referenziert): Eine Reformulierung der Physik}. Verfügbar unter: 
		
		\bibitem{pascher_derivation_2025}
		Pascher, J. (2025). \textit{Feldtheoretische Ableitung des $\beta_T$-Parameters in natürlichen Einheiten ($\hbar = c = 1$)}. Verfügbar unter: 
		
		\bibitem{pascher_qft_2025}
		Pascher, J. (2025). \textit{Vollständige Herleitung der Higgs-Masse und Wilson-Koeffizienten}. T0-Theory Project Documentation.
		
		\bibitem{pascher_units_2025}  
		Pascher, J. (2025). \textit{Natürliche Einheitensysteme: Universelle Energiekonversion und fundamentale Längenskala-Hierarchie}. Verfügbar unter: 
		
		\bibitem{katrin_2024}
		KATRIN-Kollaboration. (2024). \textit{Direkte Neutrino-Massenmessung basierend auf 259 Tagen KATRIN-Daten}. arXiv:2406.13516.
		
		\bibitem{nufit_2024}
		Esteban, I., et al. (2024). \textit{NuFit-6.0: Aktualisierte globale Analyse dreifarbiger Neutrino-Oszillationen}. J. High Energy Phys. 12, 216.
		
		\bibitem{cosmology_2024}
		Planck-Kollaboration. (2024). \textit{Planck 2024 Ergebnisse: Kosmologische Parameter und Neutrino-Massen}. Astron. Astrophys. (eingereicht).
		
		\bibitem{gell_mann_1964}
		Gell-Mann, M. (1964). \textit{A schematic model of baryons and mesons}. Physics Letters, 8(3), 214--215.
		
		\bibitem{mendeleev_1869}
		Mendeleev, D. (1869). \textit{Über die Beziehungen der Eigenschaften zu den Atomgewichten der Elemente}. Zeitschrift für Chemie, 12, 405--406.
		
		\bibitem{muon_g2_2023}
		Muon g-2 Collaboration. (2023). \textit{Measurement of the positive muon anomalous magnetic moment to 0.20 ppm}. Phys. Rev. Lett. 131, 161802.
		
	\end{thebibliography}
